\chapter{Extensionalität und Mengenbildung}

\section{Gleichheit, Teilmengen und leere Menge}

\begin{theorem}[Elementkriterium für Gleichheit und Ungleichheit]
\[
\begin{aligned}
 A=B&\ \Longleftrightarrow\
           \forall x\,(x\in A\Longleftrightarrow x\in B),\\
 A\ne B&\ \Longleftrightarrow\
 \exists x\,\bigl((x\in A\land x\notin B)
                  \lor(x\in B\land x\notin A)\bigr).
\end{aligned}
\]
\end{theorem}

\begin{theorem}[Grundgesetze der Inklusion]
\[
\begin{gathered}
 A\subseteq A,\qquad
 (A\subseteq B\land B\subseteq C)\Rightarrow A\subseteq C,\\
 A=B\ \Longleftrightarrow\ A\subseteq B\land B\subseteq A,\\
 A\subsetneq B\ \Longleftrightarrow\
 A\subseteq B\land\exists x\in B\,(x\notin A).
\end{gathered}
\]
\end{theorem}

\begin{theorem}[Eindeutigkeit und Kennzeichen der leeren Menge]
\[
 A=\varnothing\ \Longleftrightarrow\ \neg\exists x\,(x\in A),
 \qquad A\ne\varnothing\ \Longleftrightarrow\ \exists x\,(x\in A).
\]
Für jede Menge $A$ gilt $\varnothing\subseteq A$.
Ferner gilt $A\subseteq\varnothing\Longleftrightarrow A=\varnothing$.
\end{theorem}

\section{Ungeordnete Paarmengen}

\begin{definition}[Einzelmenge und Paarmenge]
Die Schreibweisen $\{a\}$ und $\{a,b\}$ bezeichnen die Mengen mit
den Elementkriterien
\[
 x\in\{a\}\Longleftrightarrow x=a,\qquad
 x\in\{a,b\}\Longleftrightarrow x=a\lor x=b.
\]
Die Elemente $a,b$ brauchen nicht verschieden zu sein.
\end{definition}

\begin{theorem}[Gleichheit von Einzel- und Paarmengen]
\[
\begin{gathered}
 \{a\}=\{b\}\Longleftrightarrow a=b,\qquad
 \{a,b\}=\{a\}\Longleftrightarrow a=b,\\
 \{a,b\}=\{c,d\}\Longleftrightarrow
       (a=c\land b=d)\lor(a=d\land b=c).
\end{gathered}
\]
\end{theorem}

\section{Aussonderung und Ersetzung}

\begin{theorem}[Eindeutigkeit und Inklusion der Aussonderung]
Die Menge $A_\varphi=\{x\in A\mid\varphi(x)\}$ existiert eindeutig.
Es gelten
\[
\begin{gathered}
 A_\varphi\subseteq A,\\
 A_\varphi=A\Longleftrightarrow\forall x\in A\,\varphi(x),\qquad
 A_\varphi=\varnothing\Longleftrightarrow
                       \forall x\in A\,\neg\varphi(x).
\end{gathered}
\]
\end{theorem}

\begin{theorem}[Vergleich ausgesonderter Mengen]
Für $A_\varphi=\{x\in A\mid\varphi(x)\}$ und
$A_\psi=\{x\in A\mid\psi(x)\}$ gelten
\[
\begin{aligned}
 A_\varphi\subseteq A_\psi
   &\Longleftrightarrow
        \forall x\in A\,(\varphi(x)\Rightarrow\psi(x)),\\
 A_\varphi=A_\psi
   &\Longleftrightarrow
        \forall x\in A\,(\varphi(x)\Longleftrightarrow\psi(x)).
\end{aligned}
\]
\end{theorem}

\begin{theorem}[Wiederholte und beschränkte Aussonderung]
\[
 \{x\in\{y\in A\mid\varphi(y)\}\mid\psi(x)\}
       =\{x\in A\mid\varphi(x)\land\psi(x)\}.
\]
Für $A\subseteq B$ gilt außerdem
\[
 \{x\in A\mid\varphi(x)\}\subseteq\{x\in B\mid\varphi(x)\}.
\]
\end{theorem}

\begin{theorem}[Eindeutig beschriebene Mengen zusammenfassen]
\label{b03guide:thm:termcollection}
Wenn zu jedem $x\in A$ genau eine Menge $y$ die Formel
$\varphi(x,y)$ erfüllt, gibt es genau eine Menge $B$ mit
\[
 y\in B\Longleftrightarrow\exists x\in A\,\varphi(x,y).
\]
Gibt es zu jedem $x\in A$ höchstens eine solche Menge $y$, besteht
dieselbe Aussage über die eindeutige Existenz von $B$, ohne daß für jedes $x$ ein
$y$ vorhanden sein muß. Für $A=\varnothing$ ist $B=\varnothing$.
\end{theorem}

\begin{definition}[Mengen durch Terme zusammenfassen]
Sei $t(x)$ ein Mengenterm mit festen Mengenparametern, der
für jedes $x\in A$ eine eindeutig bestimmte Menge bezeichnet.
Dann bezeichnet
\[
 \{t(x)\mid x\in A\}
\]
die Menge aller $y$, für die $\exists x\in A\,(y=t(x))$ gilt.
Ihre Existenz ist durch
Theorem~\ref{b03guide:thm:termcollection} gesichert.
Der Term ist eine Schreibabkürzung für eine durch eine Formel
eindeutig beschriebene Menge; eine gegebene Funktion wird nicht
vorausgesetzt.
Es wird damit keine indizierte Familie als eigenes Objekt eingeführt.
Für einen entsprechenden Mengenterm $t(a,b)$ bedeutet
\[
 \{t(a,b)\mid a\in A,\ b\in B\}
   =\bigcup\{\{t(a,b)\mid b\in B\}\mid a\in A\}.
\]
Diese Menge ist leer, wenn $A$ oder $B$ leer ist.
\end{definition}

\section{Grenzen der Mengenbildung}

\begin{theorem}[Allgemeines Mengenexistenzkriterium]
\label{b03guide:thm:setexistence}
Für jede Formel $\varphi(x)$ mit festen Mengenparametern sind
die folgenden Aussagen äquivalent:
\[
\begin{aligned}
 &\exists B\,\forall x\,(x\in B\Longleftrightarrow\varphi(x)),\\
 &\exists U\,\forall x\,(\varphi(x)\Rightarrow x\in U),\\
 &\exists U\,\forall x\,(\varphi(x)\Rightarrow x\subseteq U).
\end{aligned}
\]
Im Existenzfall ist $B$ eindeutig. Die durch $\varphi$ beschriebenen
Mengen lassen sich also genau dann zu einer Menge zusammenfassen,
wenn sie sämtlich Elemente einer gemeinsamen Menge sind,
gleichwertig: sämtlich Teilmengen einer gemeinsamen Menge.
\cite{koepke}
\end{theorem}

\begin{theorem}[Alle passenden Mengen zusammenfassen]
Für jede Formel $\varphi(x,y)$ mit festen Mengenparametern und
jede Menge $A$ sind die folgenden Aussagen äquivalent:
\[
\begin{aligned}
 &\exists B\,\forall y\,
       (y\in B\Longleftrightarrow\exists x\in A\,\varphi(x,y)),\\
 &\forall x\in A\,\exists C\,\forall y\,
       (y\in C\Longleftrightarrow\varphi(x,y)).
\end{aligned}
\]
Die Gesamtheit aller passenden $y$ bildet also genau dann eine
Menge, wenn dies für jedes einzelne $x\in A$ gilt.
Es ist keine Eindeutigkeit der passenden $y$ erforderlich.
Im Existenzfall ist $B$ eindeutig; für $A=\varnothing$ ist
$B=\varnothing$. \cite{koepke,pila}
\end{theorem}

\begin{definition}[Russell-Teilmenge einer gegebenen Menge]
Für jede Menge $A$ sei
\[
 R_A=\{x\in A\mid x\notin x\}.
\]
\end{definition}

\begin{theorem}[Russell-Aussage und fehlende Allmenge]
\label{b03guide:thm:russell}
Für jede Menge $A$ gilt
\[
 R_A\subseteq A,\qquad R_A\notin A,\qquad
 \exists x\,(x\notin A).
\]
Es gibt keine Menge, die jede Menge als Element enthält:
\[
 \neg\exists U\,\forall x\,(x\in U).
\]
Ebenso gibt es zu keiner Menge $A$ eine Menge $C$ mit
$\forall x\,(x\in C\Longleftrightarrow x\notin A)$.
Die Aussagen dieses Satzes folgen bereits ohne Fundierung.
\cite{pila}
\end{theorem}

\begin{theorem}[Keine unbeschränkte Aussonderung]
Es gibt keine Menge $R$ mit
\[
 \forall x\,(x\in R\Longleftrightarrow x\notin x).
\]
Die Schreibweise $\{x\mid\varphi(x)\}$ bezeichnet deshalb ohne
gesonderten Existenznachweis nicht notwendig eine Menge.
\end{theorem}

\begin{theorem}[Keine Menge enthält alle ihre Teilmengen]
\label{b03guide:thm:missingsubset}
Für jede Menge $A$ gilt
\[
 \exists B\,(B\subseteq A\land B\notin A),\qquad
 \mathcal P(A)\nsubseteq A.
\]
Diese Aussage folgt bereits ohne Fundierung.
\end{theorem}

\begin{theorem}[Gemeinsame Elemente definierbarer Mengen]
\label{b03guide:thm:definableintersection}
Für jede Formel $\varphi(Y)$ mit festen Mengenparametern gilt:
Wenn $\exists Y\,\varphi(Y)$, gibt es genau eine Menge $D$ mit
\[
 \forall x\,\bigl(
 x\in D\Longleftrightarrow
             \forall Y\,(\varphi(Y)\Rightarrow x\in Y)\bigr).
\]
Die Mengen $Y$, welche $\varphi$ erfüllen, müssen selbst keine
Menge bilden. Falls keine solche Menge $Y$ existiert, gibt es
kein $D$ mit dem angegebenen Elementkriterium.
\cite{kearnes}
\end{theorem}
